% ============================================================================
%  Chapter 2. Literature Review
%
%  \input file for the MSU thesis. Requires booktabs, float and a natbib-style
%  citation command in the parent preamble. Bibliography: accord/thesis/refs.bib
%
%  Positioning chapter. It reports no results of this thesis. Every method the
%  thesis applies is attributed here to the work that owns it.
% ============================================================================

\chapter{Literature Review}
\label{ch:lit}

This thesis builds nothing new. It takes a pipeline whose shape is standard in clinical natural
language processing, rule-based candidate mining followed by language model weak supervision
followed by a purpose-built neural detector, and runs at each stage the control that separates a real
finding from an artifact of the measurement. The controls are borrowed, each credited again wherever
it appears later, so what this chapter has to establish is narrower than a conventional review: not
that a gap exists in the space of methods, but that a gap exists in the space of things that have
been checked.

The review runs along six threads: the clinical task itself, attempted far less often than its
importance suggests; weak supervision and label noise, where a single language model annotating a
rare class falls outside the theory the field's tooling rests on; evaluation under low prevalence, a
literature now thirty-six years old that this thesis applies rather than extends; controls and
construct validity, the methodological backbone; pretraining loss as a predictor of downstream
utility, settled recently at a scale this thesis cannot approach; and selection and stage-by-stage
auditing of pipelines. Each thread supplies a constraint the next has to respect, and each thread's
gap is named in the chapter's final section, which states what all six leave.

\section{Contradiction and Inconsistency Detection in Clinical Text}
\label{sec:lit-contradiction}

\subsection{The clinical problem the task inherits}

Adverse drug events are among the best documented inpatient safety problems, and the classic
incidence study of \citet{bates1995ade} established both their frequency and that a substantial share
are preventable. The transition out of hospital concentrates the risk: unintended medication
discrepancies are found at a large fraction of admissions \citep{cornish2005discrepancies} and at
discharge \citep{wong2008reconciliation}, discontinuity between inpatient and outpatient care affects
a meaningful share of discharged patients \citep{moore2003discontinuity}, adverse events after
discharge are frequently drug-related and preventable \citep{forster2003adverse}, the document meant
to prevent them is itself unreliable in content and timing \citep{kripalani2007deficits}, and its
timeliness is associated with patient outcomes \citep{li2013timeliness}.

Two features of that literature shape the computational formulation: the most common discrepancy
category is omission, a relation between a statement and the absence of one rather than a
disagreement between two statements, an asymmetry that does not survive translation into a
sentence-pair classification task; and the two sides of a cross-note comparison are not
interchangeable, since a discharge medication list is an instruction that will be acted on while a
nursing note records what was observed, and that relation has a direction that carries the
consequence.

\subsection{Formulations inherited from natural language inference}

Natural language inference classifies a premise-hypothesis pair as entailment, contradiction, or
neutral, and the large general corpora \citep{bowman2015snli,williams2018mnli} made it a standard
probe of sentence-level semantic competence. MedNLI \citep{romanov2018mednli} carried the task into
the clinical domain, using clinician-written hypotheses over patient-record premises, and remains the
closest established benchmark to the problem studied here. Later clinical resources moved toward
multi-evidence and numerical reasoning over structured trial documents \citep{jullien2023nli4ct},
still within a document rather than across documents written by different roles.

Contradiction detection as a task distinct from entailment was posed by
\citet{demarneffe2008contradictions}, whose typology separates negation, numeric, antonymy and
structural contradiction and whose central observation is that contradiction is rarer and worse
served than the symmetric three-way framing implies. In the biomedical domain
\citet{alamri2016corpus} built a corpus of potentially contradictory claims from cardiovascular
abstracts, a resource about published claims rather than care documentation. Work at document scale
is more recent: \citet{hsu2021wikicontradiction} learn contradiction-aware comparisons over sentence
pairs within a Wikipedia article, and \citet{li2024contradoc} annotate self-contradictions in long
documents and show that even strong large language models are unreliable on cases requiring context.

A separate body of work shows how inference models reach high benchmark accuracy without doing the
intended inference: hypothesis-only baselines recover much of the label signal
\citep{poliak2018hypothesis,gururangan2018artifacts}, lexical-overlap heuristics survive until probed
adversarially \citep{mccoy2019hans}, the pattern is characterised as shortcut learning
\citep{geirhos2020shortcut}, and the clinical analogue appears in imaging, where models key on
site-identifying artifacts \citep{zech2018generalization}. Those studies diagnose a model that scores
well; this thesis has a pipeline that scores poorly and needs to know whether the poor score is real.

\subsection{What has actually been built for the clinical case}

Very little, and the exception is instructive. \citet{li2015discrepancy} report an end-to-end hybrid
algorithm for medication discrepancy detection at Cincinnati Children's Hospital, comparing
structured discharge prescription lists against free-text notes for 271 patients in a Complex Care
Medical Home Program. Component performance on entity detection and attribute linkage is strong, but
combining stages drops precision, recall and $F$-measure to 71.5\%, 65.2\% and 68.2\% on discrepant
medications, well below the matched-medication figures.

Three properties of that setting explain why the numbers do not transfer, and they define the problem
this thesis takes on: one side of the comparison is structured rather than extracted from prose, the
cohort is a care-management program followed deliberately and so has a raised base rate, and the
comparison is anchored, since every prescription in the list is a candidate and no retrieval step can
silently bound recall. Remove all three, as a free-text-to-free-text comparison over an unselected
critical care cohort does, and detection becomes inseparable from retrieval: a candidate generator
sits in front of the model and caps its recall absolutely.

The extraction machinery for the medication half is mature: shared tasks established the entity and
attribute layer \citep{uzuner2010medication,henry2020n2c2}, negation handling has a standard solution
\citep{chapman2001negex}, and open tooling is available \citep{neumann2019scispacy,eyre2021medspacy}.
The layer above it is not: \citet{weiskopf2013quality} find completeness, correctness and
plausibility handled unevenly, and consistency across documents about one episode is a dimension the
reuse literature flags and the extraction literature does not target. The deployment literature is
clear about the consequence: \citet{vandersijs2006overriding} report override rates of 49\% to 96\%
for low-precision drug safety alerts, and \citet{ancker2017alertfatigue} find acceptance associated
with alert-stream complexity rather than workload, so a recall figure without the corresponding
precision says nothing about clinical usefulness.

\section{Weak Supervision and Label Noise in Clinical NLP}
\label{sec:lit-weak}

\subsection{Programmatic labelling and the assumption underneath it}

Distant supervision \citep{mintz2009distant} established the pattern of generating training labels
by projecting a structured resource onto unlabelled text and tolerating the resulting noise. Data
programming \citep{ratner2016dataprogramming} and its system realisation Snorkel
\citep{ratner2017snorkel} generalised this, allowing many noisy labelling functions and fitting a
generative model over their accuracies and correlations. The clinical track record is real:
ontology-driven weak supervision matches hand-labelled training data for entity classification
\citep{fries2021trove}, and cross-modal data programming speeds model development in medical imaging
by writing rules over an auxiliary text modality \citep{dunnmon2020crossmodal};
\citet{zhang2022wssurvey} survey the paradigm.

The guarantees require multiple sources whose errors are conditionally independent given the true
label, so that disagreement is informative and aggregation recovers a better label than any single
source. With one source, aggregation is the identity function and the generative model has no signal
to fit; this is the stated premise, not a subtlety buried in the proofs. What the literature does not
say is what to do when only one source is available and hand annotation is the resource being
economised on, which is what a single instruction-tuned model creates.

\subsection{The language model as annotator}

Since 2023 the default cheap labeller has been an instruction-tuned language model, on the strength
of results reporting that such models match or exceed crowd annotators on several classification
tasks \citep{gilardi2023chatgpt}. \citet{pangakis2023validation} give the corrective: any automated
annotation procedure must be validated against human labels on the same items, since performance
varies enough across tasks that a headline accuracy from one setting licenses nothing in another.
Calibration is a second problem, since language-model confidence values are not calibrated
probabilities without measurement \citep{guo2017calibration} and are sensitive to prompt construction
and the class priors it implies \citep{zhao2021calibrate}.

Label noise has its own literature. \citet{frenay2014labelnoise} survey classification under label
noise and separate the cases that matter: noise completely at random, noise depending on the true
class, and noise depending on the features. The distinction is load-bearing, because the strongest
noise-robust guarantees assume one of the first two, whereas an annotator with a systematically
shifted prior produces a corrupted marginal rather than symmetric noise. \citet{northcutt2021labelerrors}
show that even curated benchmark test sets carry pervasive label errors whose correction can reorder
published rankings, making the reference standard itself an object needing audit.

\subsection{Surrogate labels bias what is computed from them}

The most directly relevant result is \citet{egami2023surrogates}, who consider the now-common
workflow in which a language model annotates a document set that feeds a downstream statistical
analysis, and show that direct use of surrogate labels causes substantial bias and invalid confidence
intervals even at 80\% to 90\% surrogate accuracy. Their remedy, design-based supervised learning, is
a doubly-robust estimator combining surrogate labels with a smaller set of gold-standard labels,
giving valid inference without assuming the surrogate is good.

Two things follow, and only one is usually noticed. The headline is that high annotator accuracy is
not sufficient for a valid downstream conclusion, so reporting it is not a validation step. The
second is that the correction itself requires gold labels on a subsample, which the workflow that
used a language model to avoid annotation entirely cannot supply. Their estimand is a population
quantity in a social-science analysis rather than the training signal of a detector, and no work was
found mapping the result to the case where surrogate labels are the training targets themselves. What
transfers is the shape of the warning: accuracy on the annotator's own terms says nothing about the
validity of what is built on it.

\subsection{Counting is a different problem from classifying}

An over-calling annotator is, among other things, a bad prevalence estimator, and there is a
literature about exactly that. \citet{forman2008quantifying} formalises quantification, estimating
class prevalence rather than labelling individual members, and shows the classify-and-count estimator
is biased whenever true and false positive rates differ, which is always; his adjusted-count
correction inverts the relation using estimates of those rates. \citet{gonzalez2017quantification}
review the resulting family of estimators.

The relevance to weak supervision is direct and, in clinical natural language processing, largely
unremarked. When a weak labeller's positive rate differs from the base rate, the quantification
literature says both what the discrepancy means and what correcting it requires: estimates of the
labeller's true and false positive rates on the same distribution, which require gold labels. This
correction and the Egami correction therefore fail for the same reason in the same setting, worth
noticing before running the annotator rather than after.

\section{Evaluation Under Class Imbalance and Low Prevalence}
\label{sec:lit-metrics}

\subsection{Agreement statistics depend on prevalence}

Cohen's kappa \citep{cohen1960kappa} corrects observed agreement for the agreement expected by
chance under the observed marginals, a function of those marginals in ways that surprise readers who
take it for a corrected accuracy. \citet{feinstein1990paradoxes} set out the two paradoxes: high
observed agreement can coexist with a low kappa when the marginals are badly unbalanced, and
unbalanced marginals can raise rather than lower kappa; the companion paper
\citep{cicchetti1990resolving} works through the resolution and recommends reporting the marginals
alongside the statistic. \citet{byrt1993pabak} respond with the prevalence-adjusted bias-adjusted
kappa, an affine function of observed agreement, $\mathrm{PABAK} = 2p_o - 1$, which removes the
marginal sensitivity by removing the marginal information: it reports a rescaled observed agreement
under the name of a chance-corrected coefficient, so a system evaluated by PABAK on a set with twelve
positives is scored almost entirely on its behaviour on the negatives, and whether that is right
depends on the exchange rate between a missed contradiction and a false alarm, a clinical question
the papers in this thread do not address.

\subsection{The Matthews line and its limits}

The Matthews correlation coefficient, introduced for assessing secondary structure prediction
\citep{matthews1975mcc}, is the Pearson correlation between the binary prediction and the binary
label, computed from all four confusion matrix cells. \citet{chicco2020mcc} argue it should be
preferred to accuracy and $F_1$ under imbalance because it is high only when all four cells are
favourable, and \citet{chicco2021mccbalanced} extend the argument against balanced accuracy,
bookmaker informedness and markedness; \citet{powers2011evaluation} situates it among informedness
and markedness measures, and \citet{delgado2019kappa} conclude from kappa's behaviour under skew that
it should be avoided in favour of MCC.

Balanced accuracy, the unweighted mean of sensitivity and specificity \citep{brodersen2010balanced},
is the measure most often reached for when a prevalence-blind summary is wanted, and its implicit
weighting is the least often written down. Under heavy negative skew the ROC curve is dominated by
the majority class and can look acceptable while precision-recall behaviour sits at the no-skill
line, which is why \citet{saito2015prcurve} argue the precision-recall curve should be the primary
tool at high skew; both are reported in this thesis. Two areas under the curve computed on the same
cases are correlated, so their difference cannot be tested as though the samples were independent;
\citet{delong1988comparing} give the standard nonparametric treatment, used once here to compare a
seed ensemble against a single model. \citet{schaeffer2023mirage} make the parallel argument at
scale rather than prevalence, attributing emergent abilities in large language models to nonlinear
metric choices rather than to a change in model behaviour.

\subsection{Training under skew is not measuring under skew}

The class imbalance literature is large and mostly about fitting: resampling \citep{chawla2002smote},
characterisation of when imbalance actually hurts \citep{japkowicz2002imbalance}, a broad survey
\citep{he2009imbalanced}, cost-sensitive decision rules \citep{elkan2001cost}, loss reshaping in
several forms \citep{lin2017focal,cui2019classbalanced,cao2019ldam,menon2021logitadjust}, and
positive-unlabelled formulations \citep{bekker2020pu}. Their relevance here is bounded by a single
result: \citet{vandengoorbergh2022harm} show, in clinical risk prediction, that imbalance corrections
including undersampling, oversampling and SMOTE produce severe miscalibration without improving
discrimination as measured by AUROC, and that whatever sensitivity gain they deliver is available by
shifting the threshold instead. Such corrections move the operating point along a ranking; they do
not create ranking information.

The tools that address the measurement side sit outside that literature: decision curve analysis
\citep{vickers2006dca} evaluates a model across the threshold probabilities a decision maker might
hold, and cost-sensitive learning \citep{elkan2001cost} makes the exchange rate an explicit
parameter. Reporting guidance for clinical prediction models now covers machine learning explicitly
\citep{collins2024tripodai}, and several of its items, participant flow, outcome determination,
sample size justification and prespecified performance measures, are mechanical enough to check
before a results table exists.

\section{Controls and Construct Validity in Machine Learning Evaluation}
\label{sec:lit-controls}

This section describes the methodological backbone of the thesis, and everything in it is
borrowed; laying it out at length makes the attribution unambiguous. Each subsection names
the owner, states the design, and identifies the stage of this project it is applied to.

\subsection{Negative and positive controls in observational research}

\citet{lipsitch2010negative} formalised the negative control for observational epidemiology: a
noncausal association between exposure and outcome threatens any causal claim, so one tests for it
by finding an exposure known not to cause the outcome, or an outcome known not to be caused by the
exposure, and checking whether the analysis nonetheless produces an association. An association where
none can exist is evidence that the analysis is picking up confounding, recall bias, or a defect in
procedure, though the control does not identify which. A positive control runs the same analysis in
the opposite direction, on a case where an effect is already known to exist, so the result speaks
about the instrument rather than the question asked of it; \citet{lipsitch2010negative} do not
formalise this direction, and this thesis adapts their device to build one on the same terms. The
transfer to a machine learning pipeline is direct: replace the substantive input with one that cannot
carry the claimed information, hold everything else fixed, and see whether the reported effect
survives, evidence rather than an absence of it, because the control fires and produces a number
where the theory says there should be none.

\subsection{Randomisation tests for attribution}

\citet{adebayo2018sanity} apply the same logic to saliency maps: a model randomisation test
progressively randomises trained weights and asks whether the saliency map changes, and a data
randomisation test retrains on permuted labels and asks the same; some widely used methods produce
visually plausible maps independent of both, so those maps cannot be evidence about the model. The
design generalises past saliency to any component credited with contributing something: replace what
it consumes with something that cannot contain the credited signal and see whether the contribution
persists.

\subsection{Control tasks and selectivity}

\citet{hewitt2019control} supply the design the architecture control in this thesis is built on. A
probe trained to predict a linguistic property from a neural representation may succeed because the
representation encodes the property, or merely because the probe is expressive enough to memorise the
mapping. Their control task assigns random outputs to word types, trains the same probe on that, and
defines selectivity as the gap between linguistic and control task accuracy; common probing setups
have poor selectivity, dropout does not fix this for multilayer perceptron probes, and a more
selective representation layer can win despite lower raw accuracy. \citet{belinkov2022probing}
synthesises the literature that followed: probe accuracy alone identifies neither what a
representation encodes nor whether the model uses it, an ordinal or categorical feature a component
is credited with exploiting can be replaced with an arbitrary assignment of the same shape, and if
that assignment recovers the gain, the credited feature was not the source. The comparison must be
built into the design rather than added after a result looks suspicious, because a control chosen
after the fact is a different object.

\subsection{Leakage and measurement}

\citet{kaufman2012leakage} give the formulation, detection and avoidance of leakage in data mining:
the introduction of information about the target that should not legitimately be available. Their
detection advice is behavioural, since a model performing implausibly well relative to the task's
difficulty is the primary symptom; an auxiliary objective derivable from the model's own inputs is a
textbook instance, and the diagnosis is that the target is not independent of what the model already
has.

\citet{jacobs2021measurement} bring measurement modelling from the quantitative social sciences to
computational systems, and their framing covers the whole chapter: systems operationalise
unobservable constructs through observable proxies, and the gap between construct and
operationalisation is where harms and invalid conclusions originate. They give fairness-oriented
accounts of construct reliability and validity, arguing that disagreements about fairness definitions
frequently mask disagreements about the underlying construct. The question here is not whether a
detector reaches some score but whether what is scored, a rule-mined candidate pair judged by one
annotator against heuristic-selected evidence, is a measurement of clinical contradiction.

\paragraph{What this thesis takes, and what it does not claim.}
Table~\ref{tab:lit-controls} lists the instruments and their owners. Every row is a method
applied without modification to a new object. None is a contribution of this thesis, and the
Methodology chapter repeats the attribution at the point of use.

\begin{table}[H]
\caption[Controls and their literature sources]{Controls applied in this thesis and the work that owns each one.}
\label{tab:lit-controls}
\centering
\small
\begin{tabular}{@{}p{1.9in}p{2.1in}p{1.9in}@{}}
\toprule
Instrument & Owner & Stage it is applied to \\
\midrule
Negative control & \citet{lipsitch2010negative} & Candidate generation \\
Randomisation sanity check & \citet{adebayo2018sanity} & Architecture components \\
Control task and selectivity & \citet{hewitt2019control} & Architecture components \\
Leakage diagnosis & \citet{kaufman2012leakage} & Auxiliary training objective \\
Surrogate label bias & \citet{egami2023surrogates} & Weak supervision \\
Prevalence-dependent agreement & \citet{feinstein1990paradoxes} & Reported metrics \\
Positive control, adapted & \citet{lipsitch2010negative}, negative controls & Evaluation instrument \\
Construct validity framing & \citet{jacobs2021measurement} & The task definition \\
\bottomrule
\end{tabular}
\end{table}

\section{Pretraining Loss as a Proxy for Downstream Utility}
\label{sec:lit-loss}

A pipeline that fails can always be defended on the grounds that it was too small, and the
counter-argument requires knowing whether the evaluation instrument could have registered a
real effect, which makes the relationship between a training objective and downstream
performance directly relevant. The recent literature on it is unusually clear, and two
results carry most of the weight.

\citet{lourie2025scaling} conduct a meta-analysis of published scaling results and find that
predictable scaling, small-scale results extrapolating to larger scales, occurs in a minority of
cases, 39\% of the time by their count, because small variations in experimental setup change
scaling behaviour enough to matter: the conditions under which scaling laws hold are
uncharacterised, so an extrapolation is a hypothesis rather than a prediction.
\citet{feng2026tessera} settle the narrower question of whether pretraining loss predicts
downstream utility, at a scale no thesis can match: within a Barlow Twins framework for Earth
observation foundation models, across 395 controlled training experiments and 15 downstream tasks,
the correlation between pretraining loss and downstream performance has an absolute Pearson
coefficient below 0.2, a general result about self-supervised pretraining, obtained far from
clinical text, and stronger evidence on the general proposition than anything this thesis could
produce.

The consequence is a limitation on what may be claimed, adopted rather than resisted: the Results
chapter reports a convergence experiment on the same harness for this reason, and the Introduction
states what it licenses for the clinical null rather than offering it as a contribution to the
scaling literature.

\section{Selection, Pipelines, and Disaggregated Evaluation}
\label{sec:lit-selection}

The literature on disaggregated evaluation, and what it assumes about the selection mechanism behind
a dataset, bears on one unexplained anomaly this project reports; that literature and the anomaly it
governs are both treated at length in the Supplement rather than here, since neither is a claim about
this thesis's own contradiction detector.

\subsection{Stage-by-stage auditing of pipelines}

\citet{biswas2021fairpreprocessing} already do stage-level auditing: they study fairness introduced
by data preprocessing rather than by the classifier, evaluate 37 pipelines drawn from three sources
using causal fairness measures, and show how the local fairness of individual preprocessing stages
composes into the fairness of the pipeline as a whole. The difference from the present work is the
object audited, not the idea of auditing by stage: their metric, fairness against protected
attributes, is defined at every stage, whereas the quantity here is the validity of each stage's
evidential claim, whether a candidate pool contains the pairs it should, whether a weak labeller's
marginal is usable, whether an architectural component contributes what it is credited with, and
whether the evaluation could resolve an effect. These are stage-specific and not comparable to one
another, so they cannot be composed the way a fairness metric composes.
\citet{sculley2015debt} anticipate the general hazard, that system-level defects arise at component
boundaries rather than inside any one of them, without proposing per-stage checks.

\section{What Is Left}
\label{sec:lit-gap}

The six threads leave one opening, and it is not a methodological one: each establishes what has
been built, checked or settled elsewhere, at a scale or in a subfield this thesis borrows rather than
extends, and each stops at a version of the same edge, named below one thread at a time. No
cross-note, free-text-to-free-text medication contradiction system or benchmark was found, and the
gap sits earlier than modelling, in whether a candidate pool can be assembled whose members are the
right pairs at all. No protocol was found for a single language-model annotator at low prevalence
used as the training signal itself, and every available correction requires gold labels the workflow
was built to avoid. None of the agreement-statistics papers states, for a given evaluation set, the
exchange rate a chosen measure imposes between a missed positive and an avoided false alarm
\citep{feinstein1990paradoxes}. No published work applies the instruments of
Section~\ref{sec:lit-controls} together across a pipeline's successive stages. Neither scaling paper
addresses what a positive control should look like for a clinical task with only a handful of
confirmed positives, where the question is whether a single comparison has power at all, not whether
a trend is predictable. And neither the disaggregated-evaluation nor the stage-auditing literature
supplies a per-stage check on the validity of an evidential claim in a text pipeline.

What has not been done is to take one clinical pipeline, apply the appropriate control at every
stage, and report the outcome whichever way it falls. Each control is cheap and standard in the
subfield that owns it; none was run on this pipeline before this work. The interesting output is the
joint pattern rather than any single control result: whether failures at successive stages are
independent defects or one defect propagating, and whether a labeller's own marginal already settled
what its evaluation could show.

Put in one sentence, the question this thesis asks is narrower than the pipeline that motivates it:
whether a language-model weak labeller's positive rate, set against expert-confirmed prevalence under
an honest reading of ambiguous cases, predicts what a trained detector can then show, at a sample
size that resolves only large effects. Nothing that follows claims to resolve more than that, and
nothing that follows claims a method this chapter has attributed to someone else.
